%!TeX root=../thesis.tex

\chapter{Related Work}
\label{cap:trabalhos-relacionados}

This chapter provides a contextual synthesis of literature relevant to the empirical
evaluation of Arch-Stage. Rather than presenting a standalone systematic review, it
draws on the author's earlier qualification work---which included a tertiary systematic
mapping study on software architecture for startups and a multivocal literature review
on centralization and decentralization of architectural decision-making---and
repositions those findings as narrative context for the present thesis. The chapter also
incorporates references from practitioner grey literature where relevant; such sources
are explicitly identified as such, given that grey literature carries weaker evidential
weight than peer-reviewed publications, but often captures industry trends and real
organizational experience more rapidly than academic publication cycles allow.

% ---------------------------------------------------------------------------
\section{Centralization and Decentralization of Architectural Decisions}
\label{sec:central-decentral}

A recurring theme in both academic and practitioner literature is the tension between
centralizing architectural authority in a few experienced individuals and distributing it
across development teams. Organizations pursuing fast flow and team autonomy
increasingly favor decentralized models, yet several studies report that without
explicit governance mechanisms, decentralization leads to inconsistent decisions and
architectural drift.

The author's qualification examination explored this tension through two complementary
lenses. The first chapter of the qualification presented a tertiary systematic mapping
study surveying the existing landscape of secondary studies on software architecture for
startup and scale-up organizations, identifying recurring themes in how early-stage
teams approach architecture governance under resource and time constraints. The second
chapter conducted a multivocal literature review that synthesized both academic
publications and grey literature (practitioner blogs, conference talks, and industry
reports), covering three interrelated perspectives: deriving software architecture from
requirements, the emerging socio-technical architecture trend (which situates
architectural decisions within organizational and team structures), and the tension
between centralized and decentralized decision-making.

The multivocal review surfaced a framing that the qualification termed the ``paradox of
software architecture decision-making'': industry advocates strongly for decentralized
decisions to enable team autonomy and fast flow, yet the methods and tools available
for architectural decision-making depend heavily on the tacit knowledge of experienced
architects. Decentralization thus risks transferring decision responsibility to
practitioners who lack the conceptual vocabulary or methodological support to exercise
it effectively. The qualification committee, however, found this ``paradox'' framing
ultimately unproductive as a primary research framing: the tension it described reflects
a recognized design trade-off rather than a genuine logical contradiction, and building a
thesis around resolving a paradox risked over-extending the scope of the contribution.
This thesis consequently does not employ that framing; instead, it focuses on the
concrete, tractable problem of providing tool support for one well-defined uncertainty
management technique.

The third chapter of the qualification proposed a questionnaire-based research design
to investigate how software engineering teams actually navigate architectural decisions
in distributed settings. The proposed questionnaire targeted four dimensions: demographic
and organizational context, knowledge of modern software architecture practices,
collaboration and decision-making patterns within and across teams, and specific
challenges experienced in decentralized architectural decision-making. While this
research design would have yielded broad survey data, the qualification committee
assessed the scope as too wide and the proposed SMS and multivocal review as
insufficient as primary thesis contributions without a concrete artifact or evaluation at
the center. The present thesis responds to that feedback by reorienting around the
evaluation and refinement of Arch-Stage, a concrete tool artifact, rather than a survey
of industry practice.

% ---------------------------------------------------------------------------
\section{Socio-Technical Architecture}
\label{sec:sociotechnical}

A significant body of work, both academic and practitioner-oriented, frames
architectural decision-making not merely as a technical concern but as a
\emph{socio-technical} one---shaped by the organizational structure, communication
patterns, and cognitive constraints of the teams involved. This framing is directly
relevant to the motivation for Arch-Stage, because tool-embedded scaffolding does not
only lower the technical cost of applying ArchHypo: it also changes the social dynamics
of how architectural uncertainty is surfaced and addressed within a team.

\subsection{Conway's Law and Its Implications}
\label{subsec:conways-law}

The foundational observation underlying the socio-technical framing is Conway's
Law~\cite{conway1968}: ``Any organization that designs a system (defined broadly) will
produce a design whose structure is a copy of the organization's communication
structure.'' Conway's Law implies that architectural decisions cannot be analyzed in
isolation from the organizational and team structure in which they are made. Team
boundaries become architectural boundaries; communication bottlenecks manifest as
coupling in the system design.

Skelton and Pais~\cite{skelton2019} extended this insight into an actionable framework
with the concept of the \emph{Inverse Conway Maneuver}: rather than allowing the
organization's structure to passively determine the architecture, teams should
deliberately design team topologies to produce the target architecture. In practical
terms, this means that architectural governance---including how uncertainty is managed
and who is empowered to make decisions---must be considered a design decision with
organizational as well as technical consequences.

\subsection{Socio-Technical Architecture as an Industry Trend}
\label{subsec:sociotechnical-trend}

The InfoQ Software Architecture and Design Trends Report for 2024~\cite{infoq2024trends}
identified \emph{socio-technical architecture} as a distinct ``Early Adopter'' trend,
explicitly describing it as replacing the earlier framing of ``architecture as a team
sport.'' The report notes that the trend encompasses two related ideas: on one side,
the evolving role of architects as mentors and technical leaders rather than sole
decision-makers; on the other, the growing recognition---driven in part by the
dissemination of Team Topologies~\cite{skelton2019}---that the design of a system
must consider all the people who build, support, and maintain it~\cite{infoq2024trends}.

It is important to note that InfoQ is a practitioner-oriented publication, not a
peer-reviewed academic venue. The Trends Report reflects the editorial judgment of an
experienced editorial board drawing on community submissions, conference talks, and
practitioner case studies. As grey literature, its findings carry weaker evidential
weight than systematic empirical studies; nonetheless, they are valuable as an indicator
of what concepts are gaining traction in industrial practice and what challenges
practitioners are actively grappling with~\cite{infoq2024trends}.

In a related InfoQ article, Hunter and Stojmilova~\cite{infoq2025empowering} described
a transformation journey at a software organization adopting Team Topologies and
decentralized architectural decision-making. Their account illustrates several
socio-technical dynamics that are directly relevant to this thesis: the initial disbelief
of teams that they are actually empowered to make architectural decisions, the flood of
decisions that follows once empowerment is accepted, the consequences teams face when
decisions have unforeseen cross-boundary impacts, and the eventual self-correction as
teams develop peer accountability mechanisms. All four phases are driven by social
factors at least as much as technical ones~\cite{infoq2025empowering}.

A further grey-literature source relevant to this work is the InfoQ article by
Tibbitts~\cite{infoq2025decentralized}, which argues that ``distributed architecture is
more social than technical.'' Tibbitts draws on Andrew Harmel-Law's concept of the
Architecture Advice Process~\cite{infoq2025decentralized}---in which any team can make
any architectural decision as long as it has sought advice from affected parties and
documented the decision in an ADR---as a mechanism for replacing permission-seeking
with advice-seeking. The social shift from ``who approves this?'' to ``who should I talk
to before we do this?'' is presented as a cultural unlock that enables decentralized
governance without sacrificing coherence~\cite{infoq2025decentralized}.

\subsection{Centralization vs.\ Decentralization: The Core Tension}
\label{subsec:central-vs-decentral}

The tension between centralized architectural authority and distributed team autonomy
can be framed as follows. Centralized approaches, typified by the ``ivory tower''
architect role, ensure consistency and coherence but create bottlenecks, reduce team
ownership, and can slow delivery as the central authority becomes a gating
dependency~\cite{infoq2025empowering}. Decentralized approaches empower teams and
support fast flow but risk architectural drift, duplicated effort, and decisions made
without the contextual awareness that crosses team boundaries.

Practitioner accounts such as those of Hunter and Stojmilova~\cite{infoq2025empowering}
suggest that the resolution lies not in choosing one extreme but in providing the right
scaffolding: architectural principles to guide decisions, ADRs to document them,
advisory forums to share them, and---crucially for this thesis---tools embedded in the
daily workflow to make the practice tractable for all team members, not only
experienced architects. This is precisely the gap that Arch-Stage is designed to fill.
Where existing approaches to decentralized governance focus primarily on the
\emph{documentation} and \emph{communication} of settled decisions (ADRs, advice
forums), Arch-Stage addresses the \emph{pre-decision} space: managing the uncertainty
that precedes and motivates architectural choices.

% ---------------------------------------------------------------------------
\section{Tool Support for Architectural Decision-Making}
\label{sec:tool-support}

Several tools and platforms have been proposed to support architectural
decision-making:

\begin{itemize}
  \item \textbf{ADR tools} (e.g., adr-tools, Log4brains) focus on recording and
  versioning Architecture Decision Records (ADRs). They capture the \emph{what} and
  \emph{why} of decisions but do not model uncertainty, impact, or technical plans for
  hypotheses that are not yet decisions. ADRs document settled choices; they are not
  designed to manage the pre-decision state of uncertainty that ArchHypo addresses.

  \item \textbf{Architecture analysis tools} (e.g., ATAM-based workshops, trade-off
  analysis methods) provide structured evaluation of architectural alternatives but are
  typically heavyweight, workshop-based activities rather than continuous,
  team-embedded processes. Their value lies in one-shot or periodic analysis sessions,
  not in the ongoing, iterative management of uncertainty throughout a product's
  lifecycle.

  \item \textbf{Internal Developer Portals} such as Backstage, Port, and Cortex offer
  a catalog-centric view of services and infrastructure. While they improve
  discoverability and documentation, none of these platforms natively supports
  hypothesis engineering or architectural uncertainty management. They provide the
  integration surface that Arch-Stage exploits, but leave the hypothesis engineering
  layer entirely absent.
\end{itemize}

Arch-Stage differentiates itself by being explicitly designed around the ArchHypo
lifecycle: from hypothesis elicitation through assessment, technical planning, and
resolution. Its integration with Backstage places this lifecycle within the developer's
daily tooling environment. No existing tool, to the best of the author's knowledge,
provides this combination.

% ---------------------------------------------------------------------------
\section{Empirical Studies of Architectural Practices in Industry}
\label{sec:empirical-studies}

Industry-focused empirical studies on architectural practices remain relatively scarce
compared with tool proposals and method descriptions. Among the most relevant:

\begin{itemize}
  \item The ArchHypo evaluations by Silva et al.~\cite{silva2024tse}\cite{silva2024ieeesw}
  constitute the closest antecedents to this work. They demonstrated the value of
  hypothesis engineering in two distinct settings---a web-testing startup and a
  mission-critical government budget system---but relied entirely on manual application
  with no tool support. This thesis introduces the first evaluation of ArchHypo in a
  tool-mediated setting.

  \item Studies on architectural governance in agile organizations (e.g., the work on
  SAFe and Spotify-model variants) document how companies structure roles and
  ceremonies for architecture, but rarely examine explicit uncertainty-management
  techniques.

  \item The multivocal review conducted during the author's qualification surfaced
  practitioner reports (blog posts, conference talks, white papers) describing ad-hoc
  approaches to architectural uncertainty. These reports corroborate the need for
  structured methods but, being grey literature, carry weaker evidential value.
\end{itemize}

This thesis contributes to the body of empirical evidence by evaluating a
tool-supported application of ArchHypo with multiple teams, addressing a gap in both
the ArchHypo literature (which has not yet studied tool-mediated use) and the broader
architectural-practices literature (which lacks studies of hypothesis-engineering tools
in scaled settings).
